% ---------------------------------------------------------------------------
% Free streaming: where every multipole above rank two comes from.
%
% Written 2026-10-24 by the Paper 1 stream. The supplement went from the
% source derivatives to the matching without ever constructing the solution
% being matched ONTO -- so a reader met j_l in (102) with no account of where
% it came from, and v1.0.0's entire deliverable had no section.
%
% Every equation here is read from functions/harmonics/free_streaming.py and
% its module docstring, which quotes Annals II Appendix F pp. 379-380.
% Findings cited, none minted: B-2, T-3.
% ---------------------------------------------------------------------------
\section{Free streaming, and where every $\ell\ge3$ comes from}
\label{sec:free-streaming}

\emph{The source closes at rank two and the oscillator runs at ranks zero and
one. Nothing so far produces a multipole above the quadrupole. This section is
what does, and it is also where the line-of-sight kernel comes from.}

% ===========================================================================
\subsection{What is left when the source and the scattering go}

After last scattering $\kappa'\to0$, and $\mathcal B_{A_\ell}\equiv0$ for
$\ell\ge3$ identically. What remains of the transport equation is its
streaming couplings alone --- \emph{Annals~II} Appendix~F, p.~379, in proper
time and with the \textbf{comoving} wavenumber:
\begin{equation}
  -\dot\tau_\ell=\frac{k}{a}\Big[w_\ell\,\tau_{\ell+1}-\tau_{\ell-1}\Big],
  \qquad
  w_\ell=\frac{(\ell+1)^2}{(2\ell+3)(2\ell+1)},
  \qquad \ell\ge2 .
  \label{eq:F1}
\end{equation}
Multiplying through by $\beta_\ell$ and changing to conformal time gives the
form this bundle integrates as its reference, (F.3):
\begin{equation}
  -\big(\beta_\ell\tau_\ell\big)'
  =k\left[\frac{\ell+1}{2\ell+3}\big(\beta_{\ell+1}\tau_{\ell+1}\big)
         -\frac{\ell}{2\ell-1}\big(\beta_{\ell-1}\tau_{\ell-1}\big)\right],
  \qquad \beta_\ell=\frac{\ell!}{(2\ell-1)!!} .
  \label{eq:F3}
\end{equation}

\textbf{Each multipole is coupled to its two neighbours and to nothing else.}
That is the whole content: the ladder has no closure and no termination, so
the three ranks the source supplies generate every rank above them. \textbf{It
is also the only $\ell$-coupling in the released hierarchy} --- every other
coupling in the scheme reaches $\ell\pm1$ or $\ell\pm2$ and arrives later, by
frame, which is what \textbf{F8} draws.

\paragraph{The $\ell$-weights are where an error would hide, and it would hide
at large $\ell$.} $w_\ell\to\tfrac14$ and the two bracket weights in
\eqref{eq:F3} tend to $\tfrac12$ and $\tfrac12$: every coupling in the
recursion is a ratio of linear polynomials and tends to a constant.
\textbf{So a coupling that grew like $\ell$ would be invisible at low $\ell$
and divergent at high $\ell$} --- which is exactly the regime this paper is
about, and the reason the Appendix~F chain is the acceptance spine rather than
a smoke test. Decision \textbf{S9}.

% ===========================================================================
\subsection{The $\ell=0$ row is the rank-zero source, seen from the other side}

Continued to $\ell=0$, \eqref{eq:F1} has $w_0=\tfrac13$ and reads
\begin{equation}
  \dot\tau_0=-\frac{k}{a}\,\frac{\tau_1}{3} ,
  \label{eq:F1-ell0}
\end{equation}
the monopole fed by the dipole alone. \textbf{That is the same combination the
rank-zero gravitational source is built from}, $\mathcal
B_0=-\tfrac13\mathrm D^a\tau_a$, giving $a\mathcal B_0=-\tfrac k3\tau_1$ in
mode form.

\textbf{The two are one statement, and that is why $\tau_0\equiv0$ is a
constraint rather than initial data.} $T$ is \emph{defined} as the monopole of
$T(x,e)$, so the $\ell=0$ equation does not evolve a free quantity --- it
fixes $(\ln T)^{\textstyle\cdot}$, and $\mathcal B_0$ is a functional of the
solution rather than an external forcing. \textbf{Restoring $a\mathcal B_0$ to
$C_0$ as though it were a source therefore counts it twice}, which is finding
\textbf{T-107} and is settled at \S\ref{sec:ell-zero-row}.

\emph{This connection is worth drawing because the two appear in different
sections and look like different objects.} They are the same row of the same
hierarchy.

% ===========================================================================
\subsection{The closed form, and why the kernel is a spherical Bessel function}

Appendix~F, p.~380, states the solution directly: \emph{``the solution to the
mode coefficients with respect to the time-like integration in the flat
$(K=0)$ almost-Friedmann--Lema\^itre case are spherical Bessel functions.''}
A unit monopole released at $\eta_i$ projects onto
\begin{equation}
  \tau_\ell(\eta)=j_\ell\big(k(\eta-\eta_i)\big),
  \label{eq:analytic-projection}
\end{equation}
and in this bundle the per-multipole normalisation is \textbf{fixed by the
recursion and checked rather than assumed}
(\texttt{functions/harmonics/free\_streaming.py}, \texttt{analytic\_projection}).

\textbf{This is the one fact that makes the line-of-sight solution a solution
and not an ansatz.} The kernel $j_\ell(k\Delta\eta)$ in (102) is not a window
chosen for convenience and it is not a projection imposed on the hierarchy
from outside: \textbf{it is the hierarchy's own free-streaming solution},
evaluated along the null cone. Every statement about the $k\to\ell$ map ---
that power at multipole $\ell$ comes from $k\simeq\ell/r_\theta$, that a
feature fixed in $k$ appears at a fixed $\ell$ --- is a property of
\eqref{eq:analytic-projection} and inherits whatever is true of the recursion
that produced it.

% ===========================================================================
\subsection{Three normalisations, none privileged, and one printed slip}

Appendix~F writes one recursion three ways, and the bundle integrates all
three \emph{as printed}:

\begin{center}\small
\begin{tabular}{@{}lll@{}}
\toprule
 & \textbf{written in} & \textbf{the variable} \\
\midrule
(F.1) & proper time, comoving $k$ & $\tau_\ell$ \\
(F.3) & conformal time & $\beta_\ell\tau_\ell$ \\
(F.4) & conformal time & $\alpha_\ell^{-1}\tau_\ell$, the Ma \& Bertschinger form \\
\bottomrule
\end{tabular}
\end{center}

\noindent
with $\alpha_\ell=(2\ell+1)/\beta_\ell$. \textbf{Privileging none of them is
the point}: the three differ by $\ell$-dependent factors, so an error in a
weight shows as a disagreement between normalisations that no single one can
reveal. This is the same discipline the spectrum uses when it reports (186)
and (187)--(188) separately rather than asserting they agree.

\paragraph{(F.4) as printed carries a factor its own text removes.} The
display keeps $(2\ell+1)$ on the left, while the sentence preceding it
specifies a division by $(2\ell+1)$ that was not applied; only after that
division is it the Ma \& Bertschinger form the text says it is. \textbf{The
bundle divides} --- finding \textbf{B-2}.

\paragraph{And $k$ is comoving here, which the conventions sheet does not say
of every section.} Appendix~F's $k$ is comoving and
$k_{\rm phys}=k_{\rm com}/a$; the explicit $1/a$ in \eqref{eq:F1} against a
proper-time derivative is what fixes it. The project's conventions sheet
declares $k\equiv k_{\rm phys}$ in its \emph{harmonic-recursion} section, so
the two differ by a factor $a$ \textbf{in precisely these equations}. A bare
$k$ is therefore ambiguous at exactly the point where it matters most, and
every equation in this section is written in $k_{\rm com}$.

% ===========================================================================
\subsection{What gates this, and what each gate can and cannot see}

\begin{center}\small
\begin{tabularx}{\textwidth}{@{}l >{\raggedright\arraybackslash}X l@{}}
\toprule
\textbf{gate} & \textbf{what it tests} & \textbf{result} \\
\midrule
\textbf{F1} & the recursion against \textbf{four external hierarchies}, in two
  normalisations and two phase conventions --- criterion~1 & residuals
  $\lesssim10^{-13}$ \\
\texttt{test\_appendix\_f\_chain} & that no coupling grows with $\ell$: an
  $O(\ell)$ coupling wired in early fails it at large $\ell$ & gate on S9 \\
\texttt{sign\_control} & the \emph{negative} control --- the same recursion
  with the relative sign of the two couplings flipped. \textbf{It must fail to
  reproduce the external form}, and a suite in which it passes is not testing
  the phase & negative control \\
\eqref{eq:analytic-projection} & the integrated recursion against
  $j_\ell(k(\eta-\eta_i))$ for a unit monopole, per multipole & normalisation
  checked, not assumed \\
\textbf{T-3} & that the $\eta$ grid resolves $j_\ell(k(\eta_0-\eta))$, whose
  period in $\eta$ is $\pi/k$ & raises rather than aliasing \\
\bottomrule
\end{tabularx}
\end{center}

\paragraph{What none of them sees.} All five test the recursion \emph{with the
source switched off}. They establish that free streaming is reconstructed
faithfully and say nothing about whether the right thing is streamed:
\textbf{the matching at $\eta_*$ is a separate question and is
\S\ref{sec:w-form}'s}. \textbf{F1 at $10^{-13}$ is therefore not evidence
about the spectrum}, and reading it as such would be the error the acceptance
spine exists to prevent --- a reconstruction can be exact and be matched onto
the wrong source.
